% Vietnamese translation for OI Public Library
% Source-File: IOI中国国家候选队论文/2006/黄劲松/黄劲松.doc
\documentclass[11pt]{article}
\usepackage{oipl-vn}

\newcommand{\Oh}{\mathcal{O}}

\title{Quy hoạch động tham lam}
\author{Huang Jinsong\\Trường Trung học Keqiao, huyện Thiệu Hưng, Chiết Giang}
\date{2006}

\begin{document}
\maketitle

\origtitle{Quy hoạch động tham lam: bàn về ứng dụng tư tưởng tham lam trong quy hoạch động}
\sourcepath{IOI China candidate team papers / 2006 / Huang Jinsong.doc}

\renewcommand{\abstractname}{Tóm tắt}
\begin{abstract}
Tham lam và quy hoạch động là hai kỹ thuật thường gặp trong thi tin học. Bài
viết bàn về cách đưa tư tưởng tham lam vào quá trình thiết kế quy hoạch động:
trước hết giải thích vì sao điều đó cần thiết và khả thi, sau đó trình bày hai
ứng dụng chính qua ví dụ: dùng tham lam để xác lập trạng thái, và dùng tham
lam để loại bỏ quyết định dư thừa nhằm tối ưu hóa thuật toán.
\end{abstract}

\noindent\textbf{Từ khóa:} tham lam; quy hoạch động; trạng thái; độ phức tạp
thời gian.

\section{Dẫn nhập}

Khi mô hình bài toán quá phức tạp, quy hoạch động thường gặp hai trở ngại:
trạng thái quá lớn và chuyển trạng thái khó xác định. Trong những tình huống
đó, phân tích tham lam không nhất thiết thay thế quy hoạch động, nhưng có thể
giúp ta nhìn ra tính chất ẩn của nghiệm tối ưu hoặc chỉ ra phần dư thừa trong
một công thức chuyển trạng thái. Khi ấy quy hoạch động trở nên sáng sủa hơn.

Một bài toán thích hợp với quy hoạch động thường có các đặc điểm sau:
\begin{enumerate}
  \item lời giải có thể chia thành một chuỗi giai đoạn quyết định;
  \item mỗi giai đoạn có một số trạng thái;
  \item mỗi quyết định đưa trạng thái hiện tại sang trạng thái của giai đoạn
  tiếp theo;
  \item từ một giai đoạn trở đi, dãy quyết định tối ưu không phụ thuộc vào cách
  ta đi tới trạng thái ấy;
  \item chi phí chuyển giữa các trạng thái được định nghĩa rõ ràng.
\end{enumerate}

Trong thực tế, khó khăn là nhận ra trạng thái đúng hoặc làm cho số chuyển trạng
thái đủ nhỏ. Tư tưởng tham lam hữu ích ở đúng hai điểm này: nó giúp rút ra
tính chất của nghiệm tối ưu, hoặc giúp chứng minh rằng một số quyết định chắc
chắn không cần xét.

\section{Dùng tham lam để xác lập trạng thái}

Trạng thái là trung tâm của quy hoạch động, nhưng thông tin trạng thái trong
bài toán thường bị che giấu. Nếu biết nhìn vào đặc điểm cực trị của nghiệm, ta
có thể tách được phần thông tin cần lưu.

\subsection{Ví dụ 1: Nỗi phiền của chú ếch}

Có \(n\) lá sen, \(1\le n\le 1000\), nằm tại các đỉnh của một đa giác lồi và
được đánh số theo chiều kim đồng hồ. Một chú ếch bắt đầu ở lá số \(1\), muốn
đi qua mỗi lá đúng một lần, và muốn tổng khoảng cách nhảy là nhỏ nhất. Từ một
lá, nó có thể nhảy tới bất kỳ lá nào khác.

Nhìn bề ngoài, đây là một bài TSP trên \(n\) điểm. Điểm đặc biệt là các điểm
nằm trên một đa giác lồi. Xét tư tưởng 2-opt của TSP: nếu một hành trình có
hai cạnh cắt nhau, ta có thể đổi nối hai cạnh ấy để thu được hành trình ngắn
hơn. Vì vậy trong hành trình tối ưu của bài này không có hai cạnh cắt nhau.

Từ đó suy ra bước đầu tiên từ đỉnh \(1\) chỉ có thể tới \(2\) hoặc \(n\). Nếu
đi tới một đỉnh khác, đường đi còn lại buộc phải tạo một giao cắt trong đa
giác lồi. Sau bước đầu, bài toán còn lại vẫn là bài cùng dạng trên một đoạn
liên tiếp của các đỉnh.

Đặt \(f(i,L,0)\) là độ dài nhỏ nhất khi bắt đầu từ đỉnh \(i\) và cần đi qua
các đỉnh liên tiếp
\[
  \{i,i+1,\ldots,i+L-1\},
\]
mỗi đỉnh đúng một lần. Tương tự, \(f(i,L,1)\) là độ dài nhỏ nhất khi bắt đầu
từ đỉnh \(i+L-1\) trên cùng đoạn ấy. Khi \(L=1\), cả hai giá trị bằng \(0\).
Với \(dist(a,b)\) là khoảng cách Euclid giữa hai đỉnh, ta có
\[
\begin{aligned}
f(i,L,0)=\min\{&
dist(i,i+1)+f(i+1,L-1,0),\\
&dist(i,i+L-1)+f(i+1,L-1,1)\},
\end{aligned}
\]
và
\[
\begin{aligned}
f(i,L,1)=\min\{&
dist(i+L-1,i+L-2)+f(i,L-1,1),\\
&dist(i+L-1,i)+f(i,L-1,0)\}.
\end{aligned}
\]
Đáp án là \(f(1,n,0)\), độ phức tạp \(\Oh(n^2)\).

Ở đây phân tích tham lam không trực tiếp tạo thuật toán tham lam. Nó chỉ chứng
minh tính chất ``không giao cắt'', từ đó trạng thái quy hoạch động trở nên tự
nhiên.

\subsection{Ví dụ 2: Tian Ji đua ngựa}

Tian Ji và vua Tề mỗi người có \(N\le 2000\) con ngựa. Mỗi lượt, một con ngựa
của hai bên thi với nhau; thắng được \(200\), thua mất \(200\), hòa được
\(0\). Tốc độ từng con ngựa đã biết. Cần sắp xếp ngựa của Tian Ji để tối đa hóa
tiền thắng.

Bài toán có thể mô hình hóa thành ghép cặp trọng số cực đại trong đồ thị hai
phía, nhưng cách đó quá nặng cho \(N=2000\). Ta xét các quyết định tham lam khi
ngựa của vua Tề được đưa ra từ mạnh tới yếu:
\begin{itemize}
  \item nếu ngựa mạnh nhất còn lại của Tian Ji không thắng được ngựa mạnh nhất
  của vua, nên dùng ngựa yếu nhất của Tian Ji để thua lượt ấy;
  \item nếu ngựa mạnh nhất của Tian Ji thắng được, nên dùng nó để thắng ngay;
  \item nếu hai ngựa mạnh nhất hòa nhau, có hai khả năng hợp lý: hòa bằng ngựa
  mạnh nhất, hoặc hy sinh ngựa yếu nhất.
\end{itemize}

Hai lựa chọn đầu có thể chứng minh bằng trao đổi: nếu một phương án tốt hơn
không ghép các con ngựa ấy với nhau, đổi đối thủ giữa hai cặp không làm kết
quả xấu đi. Lựa chọn thứ ba tạo phân nhánh, vì cả hòa và hy sinh đều có thể là
tối ưu trong những bộ dữ liệu khác nhau.

Điểm quan trọng là sau phân tích tham lam, ta biết Tian Ji chỉ cần lấy ngựa từ
hai đầu của dãy đã sắp theo tốc độ. Từ đó có trạng thái quy hoạch động.

Sắp ngựa hai bên giảm dần. Gọi \(score(a,b)\) là \(200,0,-200\) tùy ngựa
\(a\) thắng, hòa hay thua ngựa \(b\). Đặt \(dp(i,j)\) là lợi nhuận lớn nhất sau
\(i\) lượt, trong đó Tian Ji đã dùng \(j\) con mạnh nhất và \(i-j\) con yếu
nhất. Khi đấu với con ngựa thứ \(i\) của vua, ta hoặc lấy thêm một con mạnh,
hoặc lấy thêm một con yếu:
\[
\begin{aligned}
dp(i,j)=\max\{&
dp(i-1,j-1)+score(T_j,Q_i),\\
&dp(i-1,j)+score(T_{n-(i-j)+1},Q_i)\}.
\end{aligned}
\]
Các chỉ số không hợp lệ bị bỏ qua. Bài toán trở thành quy hoạch động
\(\Oh(n^2)\).

\section{Dùng tham lam để tối ưu quy hoạch động}

Nhiều bài có trạng thái và công thức chuyển trạng thái dễ thấy, nhưng công
thức ban đầu quá chậm. Phân tích tham lam lúc này giúp nhận ra các quyết định
không thể tối ưu.

\subsection{Ví dụ 3: Gộp đá}

Có \(n\le 1000\) đống đá trên một hàng. Mỗi lần chỉ được gộp hai đoạn kề nhau;
chi phí bằng tổng số đá của đoạn mới. Cần chi phí nhỏ nhất để gộp tất cả thành
một đống.

Quy hoạch động kinh điển đặt \(m(i,j)\) là chi phí nhỏ nhất để gộp các đống từ
\(i\) tới \(j\), và \(w(i,j)\) là tổng số đá đoạn ấy:
\[
  m(i,j)=\min_{i\le k<j}\{m(i,k)+m(k+1,j)+w(i,j)\}.
\]
Có \(\Oh(n^2)\) trạng thái và mỗi trạng thái xét \(\Oh(n)\) điểm chia, nên
tổng thời gian là \(\Oh(n^3)\).

Ta tìm phần dư thừa trong việc duyệt mọi \(k\). Gọi \(s(i,j)\) là điểm chia
đạt cực tiểu cho \(m(i,j)\). Với bài gộp đá, có thể chứng minh bằng bất đẳng
thức tứ giác rằng quyết định tối ưu đơn điệu:
\[
  s(i,j-1)\le s(i,j)\le s(i+1,j).
\]
Vì vậy khi tính \(m(i,j)\), ta chỉ cần duyệt \(k\) trong đoạn hẹp
\([s(i,j-1),s(i+1,j)]\), thay vì toàn bộ \([i,j-1]\). Tổng thời gian giảm xuống
\(\Oh(n^2)\).

Đây là một biểu hiện điển hình của tư tưởng tham lam trong tối ưu quy hoạch
động: nếu một quyết định nằm ngoài khoảng có thể tối ưu, ta không xét nó nữa.

\subsection{Ví dụ 4: The Lost House}

Một con ốc sên mất nhà tại một lá của một cây. Nó bắt đầu từ gốc. Một số nút
trong có côn trùng sống ở đó; nếu đi tới nút ấy, côn trùng cho biết nhà có nằm
trong cây con của nút đó hay không. Nếu không có côn trùng, ốc sên phải tự
quyết định thứ tự thăm các cây con. Mỗi cạnh đi qua tốn \(1\), và xác suất nhà
nằm ở mỗi lá là như nhau. Cần kỳ vọng số bước nhỏ nhất.

Đặt:
\begin{itemize}
  \item \(F_a(u)\): tổng kỳ vọng đóng góp khi nhà nằm trong cây con của \(u\);
  \item \(F_b(u)\): thời gian để duyệt xong cây con của \(u\) khi nhà không nằm
  trong đó;
  \item \(Leaves(u)\): số lá trong cây con của \(u\).
\end{itemize}
Đáp án tại gốc là
\[
  \frac{F_a(root)}{Leaves(root)}.
\]

Giả sử \(u\) có các con được thăm theo thứ tự
\[
  S_1,S_2,\ldots,S_k.
\]
Khi tính theo thứ tự ấy, ta có thể dùng đoạn giả mã:
\begin{verbatim}
Fa[u] = 0; Fb[u] = 0;
for i = 1 to k do begin
  Fa[u] = Fa[u] + (Fb[u] + 1) * Leaves[S[i]] + Fa[S[i]];
  Fb[u] = Fb[u] + Fb[S[i]] + 2;
end;
\end{verbatim}

Vấn đề còn lại là chọn thứ tự các con. Duyệt hết mọi hoán vị tốn
\(\Oh(nk!)\); quy hoạch động trạng thái tập con cũng chỉ tạm chấp nhận được
khi \(k\le 8\), nhưng không phải trọng tâm của bài.

Ta xét hai con kề nhau \(x,y\). Giả sử trước cặp này đã tốn \(P\) bước. Nếu
thăm \(x\) rồi \(y\), phần phụ thuộc vào hai con là
\[
  (P+1)Leaves(x)+F_a(x)+(P+F_b(x)+3)Leaves(y)+F_a(y).
\]
Nếu thăm \(y\) rồi \(x\), biểu thức tương tự thu được bằng cách đổi vai \(x\)
và \(y\). Lấy hiệu hai biểu thức, các hạng \(P,F_a(x),F_a(y)\) triệt tiêu. Do
đó \(x\) nên đứng trước \(y\) khi
\[
  (F_b(x)+2)Leaves(y)\le (F_b(y)+2)Leaves(x),
\]
hay tương đương
\[
  \frac{F_b(x)+2}{Leaves(x)}
  \le
  \frac{F_b(y)+2}{Leaves(y)}.
\]

Quan hệ so sánh này có tính bắc cầu, nên ta chỉ cần sắp xếp các con theo tỷ
số trên rồi áp dụng công thức cây. Chi phí tại mỗi nút là chi phí sắp các con;
tổng quát có thể viết là \(\Oh(nk\log k)\) nếu \(k\) chặn bậc lớn nhất. So với
duyệt hoán vị, đây là một cải tiến rất lớn.

\section{Tổng kết}

Bốn ví dụ trên minh họa hai vai trò của tham lam trong quy hoạch động.
\begin{itemize}
  \item Khi trạng thái bị che giấu, tham lam giúp rút ra tính chất cấu trúc của
  nghiệm tối ưu, như đường đi không giao cắt trong đa giác lồi hoặc việc chỉ
  lấy ngựa từ hai đầu dãy đã sắp.
  \item Khi trạng thái đã rõ nhưng thuật toán chậm, tham lam giúp loại bỏ quyết
  định dư thừa, như quyết định đơn điệu trong gộp đá hoặc thứ tự thăm cây con
  trong \textit{The Lost House}.
\end{itemize}

``Quy hoạch động tham lam'' không phải một thuật toán riêng biệt. Nó là một
cách phân tích: hiểu sâu bài toán để tìm thông tin ẩn, rồi dùng thông tin ấy
để làm quy hoạch động gọn hơn hoặc nhanh hơn.

\section*{Lời cảm ơn}

Tác giả cảm ơn toàn thể đội tuyển Chiết Giang NOI 2005 đã hỗ trợ, và thầy Wu
Jianfeng của Trường Trung học Keqiao vì hướng dẫn kỹ thuật.

\begin{thebibliography}{9}
\bibitem{liurujia}
Liu Rujia và Huang Liang, \textit{Nghệ thuật thuật toán và thi tin học}.

\bibitem{heuristics}
Zbigniew Michalewicz và David B. Fogel, \textit{How to Solve It: Modern
Heuristics}.

\bibitem{dpopt}
Mao Ziqing, \textit{Kỹ thuật tối ưu hóa thuật toán quy hoạch động}, bài viết
đội tuyển IOI 2001.
\end{thebibliography}

\section*{Ghi chú về nguồn}

Tệp Word gốc dùng nhiều công thức nhúng và đối tượng hình cũ. Một số công thức
được khôi phục từ ngữ cảnh thuật toán; hình minh họa và phụ lục chương trình
dài không được chép lại. Nội dung dịch tập trung vào lập luận thuật toán chính
và các công thức cần thiết.

\end{document}
